\chapter{Literature Review}

    \paragraph{} Our literature review began with a deliberate return to first principles, a slow and careful journey through the foundations of the methods we hoped to employ, because we believed that only by understanding these approaches at their roots could we appreciate both the strength they offer and the constraints they inevitably carry. We immersed ourselves in the conceptual architecture of knowledge graphs and the evolving capabilities of large language models, allowing ourselves to see how these technologies rise from their theoretical underpinnings and extend into practical application, while also revealing the subtle tensions between structured representation and probabilistic reasoning. Through this grounding, the benefits of these methods came into sharper focus, yet so too did their shortcomings, which encouraged us to think critically about the kinds of challenges we might face and the kinds of risk mitigation strategies we would eventually need to design.
    
    \paragraph{} With that foundation firmly in place, we turned to the broader landscape of real-world applications, studying how these techniques have been adopted across a variety of disciplines and observing the different types of problems they were intended to solve, as well as the new questions they occasionally introduced. Because we were now more attuned to the conceptual nuances, we began reading each paper with a more discerning perspective, attentive not only to what worked but also to what seemed fragile, incomplete or unsuitable for certain contexts, and this growing critical awareness allowed us to map not just the successes but also the limitations of the field.

    \paragraph{} To avoid the narrowness that comes from focusing solely on our own domain, we intentionally sought out examples from industries far beyond automotive engineering, exploring how knowledge graphs and related techniques have been used in pharmaceuticals, metal manufacturing and various compliance-driven environments. As we compared these diverse applications, we began to notice recurring patterns in how these methods succeeded, how they faltered and how they shaped the organisations that adopted them, which helped us develop a more balanced and objective understanding of their true capabilities. This cross-industry perspective encouraged us to think not only about what these methods might offer in our own case study but also about how our framework might contribute more generally to the evolving conversation around knowledge-based reasoning systems.

    \newpage
    \paragraph{} Only after we had absorbed both the conceptual and cross-domain insights did we begin exploring the more technical pathways that support these approaches, investigating how different implementation techniques influence the overall performance and adaptability of such systems. This phase led us into the details of neural networks, recommender system architecture and emergent tools such as the model context protocol, each of which offered distinct possibilities for how we might realise our own design vision. The variety of technological routes available to us made it clear that there was no single path to the solution we sought, but rather a constellation of approaches, each with its own advantages and trade-offs.
    
    \paragraph{} Given the pace at which knowledge graph research continues to expand, with new contributions appearing at a remarkable rate, we made a conscious effort to prioritise the most current, relevant and impactful studies, ensuring that our understanding remained tightly connected to the state of the art. By constantly revisiting new research, refining our interpretations and considering how each development might influence our own implementation strategy, we were able to shape a literature review that was neither static nor superficial, but instead reflective, evolving and deeply informed. Through this process, our review became not just a survey of what others had done, but a narrative of learning that equipped us with the clarity and judgement needed to construct the framework at the heart of our thesis.

\newpage
\section{Foundational Studies \& Core Technologies}

    \subsection{Knowledge Graphs}

        \paragraph{} Knowledge Graphs (KGs) are structured representations of knowledge, typically modelling entities (nodes) and their semantic relationships (edges). The concept gained prominence in 2012 when Google introduced the term \textit{Knowledge Graph} to enhance search results by organising facts about real-world entities. However, the roots of KGs trace back to earlier ideas in AI and the Semantic Web, such as semantic networks and ontologies for knowledge representation [1,7,9]. A KG can be seen as a directed graph of triples (head entity, relation, tail entity. By capturing factual connections in a graph format, KGs support intuitive querying and integration of heterogeneous data sources. KGs have become foundational in many applications, including question answering, semantic search, recommendation, and language understanding [2,9,10].

        \paragraph{} A key step in building a knowledge graph is \textbf{ontology engineering} - designing a schema or formal ontology that defines the types of entities and relationships in the domain [3,5]. Ontologies provide a shared vocabulary and rigorous definitions of concepts, essentially serving as the blueprint for a KG’s structure. With an ontology in place, one can perform \textbf{knowledge extraction} to populate the graph. This typically involves extracting entities, attributes, and relations from unstructured or semi-structured sources (e.g. text documents, databases) using Natural Language Processing (NLP) techniques [4,11]. For example, named entity recognition and relation extraction methods can turn sentences or tables into candidate triples for the KG. As knowledge is gathered from multiple sources, \textbf{knowledge fusion} becomes critical - i.e. merging and aligning information about the same real-world entities from different data silos [5,6]. Knowledge fusion approaches often rely on entity resolution or ontology alignment to integrate new facts into a unified graph. Once constructed, a KG may still be incomplete; \textbf{knowledge extension} refers to adding new entities or relations to broaden the graph’s scope, and \textbf{knowledge graph completion} (or \textbf{knowledge reasoning}) refers to inferring or predicting missing relations between known entities [8,12,16]. These processes benefit from both logical reasoning and machine learning, ensuring the KG grows in coverage and accuracy over time. Maintaining a KG also entails knowledge updating – revising or removing facts as the underlying information changes. This is especially important in dynamic domains where facts evolve [7,14].

        \paragraph{Static Vs. Dynamic Knowledge Graphs:}
        Traditionally, most KGs have been static, meaning they represent a snapshot of knowledge that changes only through occasional manual updates. A static KG provides a stable, consistent view of data, but one that can become outdated if knowledge in the world progresses. In contrast, dynamic knowledge graphs are designed to continuously evolve by integrating new data and forming new connections in real time [11,12]. The core difference is adaptability: a dynamic KG supports ongoing knowledge extension and updating, reflecting changes such as emerging entities, new facts, or temporal events. Dynamic updates focus on extracting and integrating new information from multi-source data streams, thus enabling the continuous evolution of the KG [12,18]. This adaptive capability is crucial for domains like news, finance, or social data, where facts rapidly change. Dynamic KGs can be seen as “living” knowledge bases that grow and adjust as new inputs arrive, whereas static KGs are fixed after their construction.

        \paragraph{}
        Recent literature has also distinguished \textbf{temporal knowledge graphs} as a related concept: these are graphs in which each triple is augmented with a timestamp to capture temporally scoped facts (forming quadruples) [18,19]. A temporal KG explicitly represents time-evolving relationships, allowing queries about what was true when. Similarly, researchers discuss \textbf{event knowledge graphs} that focus on nodes representing events (with time and participants) as first-class entities [18]. Managing temporal and dynamic KGs presents additional challenges in storing versions of facts over time and ensuring consistency as the graph expands or changes.

        \paragraph{}
        Beyond text-based data, there is growing interest in \textbf{multi-modal knowledge graphs} (MKGs), which integrate information from multiple modalities (e.g. images, audio, video) alongside traditional textual knowledge. The majority of prominent KGs have historically concentrated on textual and structured data, paying comparatively little attention to visual or other sensory data. MKGs aim to bridge this gap by linking entities not only to textual descriptions but also to relevant images, videos, or even signals. For instance, an entity like Eiffel Tower in a multi-modal KG might be connected to photographs, architectural blueprints, or audio commentary, enriching the graph’s context. Incorporating these diverse data types can improve machine perception and reasoning in real-world scenarios. Research surveys on MKGs outline methods for multi-modal entity and relation extraction (e.g. detecting objects in images and linking them to entities) and for multi-modal representation learning that fuses features from text and vision [9,27,28]. The inclusion of multiple modalities poses new theoretical questions on how to align and merge different feature spaces, but it ultimately leads to more comprehensive knowledge representations of the world.

        \paragraph{}
        Underpinning the development of KGs are the platforms and technologies that store and query graph-structured data. Two prevalent paradigms are the Resource Description Framework (RDF) with SPARQL for triple stores and labelled property graph databases with graph query languages (like cypher). A notable example is \textbf{neo4j}, a popular graph database that, while not a knowledge graph itself, provides the infrastructure for creating and managing KGs under the labelled property graph model [14]. Neo4j has been instrumental in the practical growth of KG technology by offering efficient storage of nodes and relationships, indexing, and traversal algorithms that can scale to large graphs. It supports ontology-like schemas through labels and relationship types and is often used in enterprise settings to implement knowledge-driven solutions. In summary, knowledge graphs marry principles from semantic technologies and graph theory to create a foundation for intelligent systems. Through ontology engineering, extraction, fusion, and iterative updates, KGs provide a dynamic yet structured \textbf{knowledge base} that can be leveraged by various AI applications [1,9,20].

    \newpage
    \subsection{Natural Language Processing \& Large Language Models}

        \paragraph{}
        There is a strong symbiosis between knowledge graphs and Natural Language Processing (NLP) [1,2] On one hand, NLP techniques enable the \textbf{construction of KGs} from unstructured text by extracting entities, relations, and facts 3. Classical IE (\textbf{information extraction}) pipelines use methods like Named Entity Recognition (to identify key entities in text), relation extraction (to detect semantic relations between entities in a sentence), and coreference resolution (to link mentions of the same entity) [3,4]. These tools can populate a KG from corpora such as news articles or the web. For example, an NLP system might read that “Alan Turing Wrote On Computable Numbers In 1936” and add a triple \textit{(Alan Turing, authorOf, On Computable Numbers)} with a date attribute to a knowledge graph 1. Such \textbf{knowledge extraction} from text is an essential step to keep KGs rich and up-to-date [4,5]. Conversely, knowledge graphs can aid NLP tasks by providing structured world knowledge that helps disambiguate text or answer complex queries 6. Early QA systems, for instance, used KGs to answer natural language questions by translating them into graph queries (SPARQL or path queries) 7. The KG acts as a semantic memory that an NLP system can consult for facts, a pattern that has been particularly evident in \textbf{knowledge-based question answering} (KG-QA) [7,8].
                \paragraph{}
        In recent years, \textbf{Large Language Models (LLMs)} have revolutionised NLP by learning from massive text corpora [9]. Models such as GPT-4 or PaLM possess an impressive ability to generate and understand text, but they primarily capture knowledge in implicit, probabilistic ways. Integrating LLMs with explicit knowledge graphs is a burgeoning area of research, as each offers complementary strengths [10] KGs provide \textbf{symbolic, curated knowledge} that is exact and transparent (e.g., a list of country-capital pairs), while LLMs provide \textbf{statistical language knowledge} and generalisation ability (e.g., the ability to infer that a capital is a kind of city and answer in full sentences) . By fusing these, one can potentially address the limitations of both technologies  One integration strategy is KG-enhanced LLMs, where a language model is augmented with information retrieved from a KG to improve its factual accuracy and reduce \textit{“hallucinations”} [12,13]. For example, an LLM can be hooked up to a KG-based retrieval system that supplies relevant triples when the model is answering a question, ensuring the answer aligns with the KG’s facts. This approach, often termed \textbf{retrieval-augmented generation} (RAG), has shown promise in making LLM outputs more trustworthy by grounding them in a knowledge graph [13]. The reverse strategy is LLM-assisted KGs, where LLMs help build or refine a knowledge graph [15]. Here, the model’s vast pretrained knowledge can be used for tasks like suggesting new facts, validating the consistency of existing triples, or even filling in gaps (a form of knowledge graph completion). For instance, an LLM might be prompted with a subject and relation to predict a likely object, effectively proposing new edges for the graph [16]. Recent work demonstrates using GPT-style models to validate KG content: e.g., flagging potentially incorrect triples by cross-checking them against information implicitly known to the model or against textual evidence [17]. This LLM-enhanced KG paradigm treats the language model as a powerful extractor and verifier that can operate in tandem with the structured graph [17,18].
        \newpage
        \paragraph{}
        While the synergy of KGs and LLMs is very promising, it also poses theoretical challenges [19]A fundamental issue is the \textbf{mismatch in knowledge representation}: KGs store explicit facts in a symbolic form, whereas LLMs encode knowledge in distributed vector representations learned from text [19,20]. This can lead to representational conflicts, for example, an LLM might consider contextual nuances that are not captured in a KG’s rigid schema, or it might lack consistency in how it refers to entities across different contexts . Combining the two paradigms requires careful design to avoid introducing noise or inconsistencies [21]. Researchers have noted that naive augmentation of LLMs with KG data can sometimes inject semantic noise or bias if not done properly [22]. Another challenge is ensuring \textbf{alignment}: an LLM might use different names or formats for entities than those in the KG (the entity linking problem), requiring robust linking mechanisms . Despite these hurdles, the research trend is clearly toward tighter integration [24]. Future directions envision a \textit{neuro-symbolic loop} where LLMs and KGs iteratively refine each other, for example, the LLM populates a KG with new insights, and queries over the KG inform the LLM’s context, establishing a feedback cycle of continuous learning and updating [24,25]. This could combine the generalisation of statistical models with the precision of symbolic ones . In summary, NLP (from rule-based methods to LLMs) plays a dual role in the knowledge graph eco-system: both as a source of new knowledge (extracting structured facts from text) and as a consumer of the graph’s knowledge (improving language understanding and generation) [1,4,9]. The interplay between large language models and knowledge graphs is a cutting-edge research frontier, aiming to build AI systems that are both knowledgeable and linguistically fluent [25,26].
        

    \subsection{Reinforcement Learning}

        \paragraph{} Reinforcement Learning (RL) has emerged as an effective paradigm for certain knowledge graph tasks that can be framed as sequential decision-making problems [22,15]. In a typical RL setup, an agent interacts with an environment through a sequence of actions, receiving rewards that guide it to learn an optimal strategy. When the environment is a knowledge graph, the agent’s goal can be to navigate or modify the graph in a way that accomplishes a reasoning task. One prominent application is \textbf{knowledge graph reasoning via multi-hop traversal}. Here, the problem might be: given a query (head entity $h$ and relation $r$), find the correct tail entity $t$. Rather than one-step inference, the agent incrementally \textit{“walks”} from the head entity through a chain of relations to arrive at the answer, essentially discovering a path $h \rightarrow \dots \rightarrow t$ that satisfies the query. This can be modelled as a Markov Decision Process on the graph, where states correspond to the current node (entity) the agent is on, and actions correspond to choosing a relation and outgoing edge to follow next. The agent is rewarded when it reaches a target entity that answers the query or when it finds a valid logical path, and it is trained through repeated trials to navigate efficiently. Such RL-based approaches to \textbf{multi-hop question answering} or \textbf{link prediction} have been demonstrated in works like DeepPath and MINERVA [16,21], which use policy gradients to learn navigation policies over KGs. The appeal of RL here is that it can learn to exploit the KG’s structure and glean implicit logical rules by trial-and-error, without requiring explicit enumeration of all paths. Moreover, the paths found by the agent can serve as explanations for the inferred answer (e.g. a chain of reasoning connecting a movie to a genre via an actor).

        \newpage
        \paragraph{} Another area where RL intersects with KGs is in\textbf{ dynamic knowledge graph updating and optimisation}. In scenarios where a knowledge graph needs to be built or refined iteratively, an RL agent could decide what new information to incorporate or which sources to trust, treating the growth of the KG as a sequential decision process [11]. For example, an agent could be tasked with autonomously curating a KG by crawling sources and adding triples, receiving positive rewards for adding correct facts and negative rewards for adding incorrect ones. Over time, the agent would learn a policy that maximises the quality of the KG. This concept of using RL for \textbf{knowledge acquisition strategies} is still emerging, but it aligns with viewing KG construction as an incremental, trial-and-error process [21].

        \paragraph{} Beyond KG reasoning and construction, reinforcement learning has been applied to \textbf{knowledge-aware recommender systems} and other decision-making systems that leverage KGs [27]. In explainable recommendation, for instance, an agent might walk in a user-item knowledge graph to find a connecting path that justifies a recommendation, learning policies that balance recommendation accuracy with path length or diversity. In such cases, the KG provides the state space of possible moves linking users to items via attributes (like user → liked genre → item), and RL can help in exploring this space for optimal recommendations or explanations. In summary, RL contributes a powerful toolkit for exploring and exploiting knowledge graphs. It enables an agent to learn how to traverse complex networks of facts or how to update them, rather than relying solely on deterministic graph algorithms [22]. A recent survey by Huo (2024) provides a taxonomy of RL-based knowledge graph models and highlights diverse applications ranging from question answering to dialogue systems [21]. The survey underscores that while RL can greatly enhance the reasoning capabilities and interactivity of knowledge graphs, challenges remain in terms of efficiency (the state space of a large KG is enormous) and reliability of the learned policies. Nonetheless, reinforcement learning represents a promising avenue to imbue knowledge graphs with a degree of autonomy and adaptive reasoning, complementing the more static inference methods traditionally used in knowledge bases.

    \subsection{Recommender Systems}

        \paragraph{} Recommender systems aim to suggest items (products, content, and more.) to users based on their preferences, and knowledge graphs have been increasingly leveraged to improve recommendation quality. The integration of KGs into recommender systems gives rise to \textbf{knowledge-aware recommendation}, where the structured knowledge about users, items, and their attributes or relationships is used to enhance the recommendation process. One major benefit of using a KG is the alleviation of data sparsity and cold-start problems that plague traditional recommender systems. In cold-start scenarios (new user or new item with little interaction data), a knowledge graph can provide side information – for example, even if a new movie has no ratings yet, the KG might reveal it shares actors or genres with movies the user liked, allowing the system to infer interest. Because KGs encode high-order connections (user → liked book → same author → other books by author, etc.), they enable the system to make \textbf{diverse and serendipitous recommendations} by traversing those links [18, 19, 20]. Recent surveys note that knowledge graph-based recommendations have gained traction precisely because they can utilise this rich relational data to overcome the sparse interaction data issue.

        \newpage
        \paragraph{} Another key contribution of KGs to recommender systems is improved \textbf{explainability} of recommendations. Traditional collaborative filtering often operates as a black box, offering little insight into why an item was recommended. In contrast, with an underlying knowledge graph, one can often trace a path or a set of features that led to a recommendation. For instance, a system might recommend a film to a user and explain it by citing that \textit{“You might like this film because it is directed by Christopher Nolan, whose movies you have rated highly”}, if the KG has relations connecting the user’s liked movies to Nolan as director and then to the recommended film [19]. Such reasoning is facilitated by the graph structure. Researchers have proposed ranking paths or nodes in the KG as a way to generate human-understandable explanations. In fact, ranking entities connected to an item (like the genre, cast, or author) can highlight which features most strongly match the user’s profile, serving as an explanation. Studies have shown that providing these KG-based explanations increases user trust and transparency in the system’s outputs. This aligns with a broader trend in recommender systems toward explainable AI, where justifications are as valued as the accuracy of the recommendations.

        \paragraph{} Knowledge graphs also allow for \textbf{content-based} and \textbf{knowledge-based recommendations} that go beyond user behaviour patterns. By infusing domain knowledge, a recommender can reason about item similarity or user preferences on a conceptual level. For example, a music recommendation engine could use a music knowledge graph to understand that a user’s favourite songs are all of the subgenre \textit{“bossa nova”} and then recommend other artists classified under that subgenre, even if the user hasn’t listened to them before [19, 20]. Moreover, graph algorithms (such as random walk or graph neural networks) applied on the knowledge graph can learn low-dimensional embeddings of items that capture not only co-rating patterns but also semantic relationships from the KG. These \textbf{knowledge graph embeddings} can be used within neural recommendation models to improve predictions. One representative approach is to perform a graph convolution over the KG to aggregate features for an item from their neighbours (e.g., an artist node will aggregate information from the genres, albums, and similar artists connected to it), yielding a knowledge-enriched representation that is fed into the recommendation model.

        \paragraph{} In summary, the incorporation of knowledge graphs has expanded the capabilities of recommender systems by providing richer input features and logical connections. It mitigates common problems like sparsity, and importantly, it enables \textbf{richer user experiences} through explainable recommendations. The research community has produced comprehensive surveys in recent years charting the landscape of KG-based recommender s - covering techniques from path-based recommendation to embedding-based hybrid models, to graph neural network approaches [18, 19, 20]. These works conclude that knowledge graphs serve as an additional \textit{“semantic lens”} through which user-item interactions can be interpreted, leading to more robust and interpretable recommendation outcomes. As knowledge graphs become more prevalent and comprehensive (e.g., encompassing multimedia content or dynamic user context), we can expect recommender systems to become even more context-aware, transparent, and accurate in aligning suggestions with user needs and interests [20].

    \newpage
    \subsection{Neural Networks}
    
        \paragraph{} Neural network techniques have profoundly influenced how knowledge graphs are represented and utilised, giving rise to what some call neural knowledge graphs - a convergence of symbolic graph structures with sub-symbolic neural representations. One of the foundational developments at this intersection is \textbf{knowledge graph embedding} (KGE). The idea behind KGE is to encode entities and relations of a KG into a continuous vector space while preserving the graph’s structural and semantic information. By representing nodes and links as vectors (or matrices), we can apply the machinery of machine learning and deep learning to perform tasks like link prediction, clustering, or similarity search on the KG. Early KGE models, such as TransE, TransH, or DistMult, used shallow neural architectures or simple translational principles to embed triples: for example, TransE treats a relation as a vector translation from a head entity vector to a tail entity vector. Subsequent approaches grew more sophisticated, incorporating neural network modules - convolutional layers (as in ConvE), or full multi-layer perceptrons - to capture patterns in triples. A survey of embedding methods categorises them into translational, tensor factorisation, and n\textbf{eural network-based models} [27, 28]. Despite their differences, all these models are trained by iterating over graph triples and adjusting vector representations to score true triples highly and false triples low. The result is a set of embeddings that implicitly contain the knowledge from the graph. These embeddings can then be fed into downstream neural models (for recommendation, question answering, and more) or used to directly infer new facts by simple vector algebra. The KGE approach has been a theoretical breakthrough in that it marries the discrete world of graph semantics with the continuous world of neural optimisation, enabling gradient-based learning on knowledge graphs.
        
        \paragraph{} Building on the idea of embeddings, \textbf{graph neural networks} (GNNs) have become a key technology for deep learning on knowledge graphs. Traditional neural networks (like CNNs or RNNs) assume data in grid or sequence form, but GNNs generalise learning to arbitrary graph structures. In the context of KGs, a prominent example is the \textbf{relational graph convolutional network} (R-GCN). An R-GCN extends the concept of graph convolution to multi-relational graphs (i.e. KGs with labelled edges). It learns to aggregate information from a node’s neighbours, while accounting for the type of each edge, to update the node’s representation. For instance, given a medical knowledge graph, an R-GCN can compute an embedding for a patient node by combining information from connected nodes like diseases, treatments, and test results, using separate weight matrices for each relation type. Through multiple layers, the network propagates and mixes features along the graph’s topology, resulting in rich \textbf{knowledge-aware node embeddings}. These can be used for tasks such as predicting a missing relation (link prediction) or classifying nodes (e.g., inferring a category of an entity) [27, 28]. GNNs provide a flexible framework to incorporate both graph structure and raw features of entities into neural learning. They have achieved state-of-the-art results in many KG tasks by effectively performing \textbf{message passing} over the knowledge graph. A benefit of GNN models is that they can operate in an inductive mode - potentially generalising to new entities not seen during training by aggregating neighbour information - which is useful for dynamic or growing graphs.

        \newpage
        \paragraph{} The combination of deep neural networks with knowledge graphs also paves the way for \textbf{neuro-symbolic reasoning}. A neural knowledge graph approach seeks to blend the strengths of symbolic reasoning (which KGs excel at, providing interpretability and logical rigour) with the pattern recognition and generalisation power of neural networks. This has led to systems where neural components learn to reason over the graph. For example, reinforcement learning agents with neural policies (as discussed earlier) can learn to traverse the graph, effectively learning symbolic rules in a latent form. Another example is neural link predictors that learn to approximate logical inference (like differentiable versions of logical operators using neural attention mechanisms). There is active research in using \textbf{transformer-based models} that take relational facts as input (serialising sub-graphs) to perform reasoning, essentially treating knowledge graph completion as a sequence-to-sequence prediction problem [27, 28]. Large language models can even be fine-tuned to emit triples or walk a graph if given the right architecture and training data, which starts to blur the line between a parametric LLM and an explicit KG.
        
        \paragraph{} In essence, neural networks provide the tools to \textbf{learn representations and patterns} from knowledge graphs automatically, rather than relying solely on manual feature engineering or rule-based reasoning. Techniques like KG embeddings and graph neural networks enable capturing complex interactions in the data (such as transitive or compositional relations) via latent features. This not only improves performance on classical tasks like link prediction and entity classification, but also opens new possibilities: e.g., combining visual or textual encoders with graph encoders for multi-modal KGs, or using generative models to synthesis new graph fragments. The era of neural networks and deep learning has thus transformed knowledge graph research, yielding approaches where learning and reasoning go together [27, 28]. As Fanghua Yu (2024) notes, we are entering the \textit{“era of the neural knowledge graph,”} aiming to merge interpretable symbolic structures with powerful deep learning, in pursuit of AI systems that can both learn from raw data and reason with structured knowledge. The ongoing challenge is to maintain the balance between these two: preserving the explicit meaning and reliability of knowledge graphs while leveraging neural networks to scale up and enrich what can be inferred from them. With continued research, the integration of KGs and neural networks is expected to produce increasingly intelligent and explainable knowledge-driven systems.

\section{Specialised Domains \& Applied Research}

    \subsection{Ontologies For Risk \& Regulatory Compliance}

        \paragraph{} Regulatory compliance and risk management often involve vast bodies of rules and standards that are complex, interdependent, and frequently changing [3,4,5,6,7,8]. Organisations must track legal requirements (financial regulations, safety codes, data privacy laws, etc.) and assess risks of non-compliance, which is a knowledge-intensive task [1,2]. To handle this complexity, many approaches turn to ontologies – formal domain models of concepts and relationships – as a foundation for building knowledge graphs that encode regulatory knowledge [3,4,5]. Ontologies provide a structured, unambiguous vocabulary for the domain, ensuring that all systems refer to entities and rules in a consistent way. This semantic consistency is crucial in compliance, where misinterpreting a term or relationship can lead to errors or legal penalties [6].

        \paragraph{} In fact, the Financial Industry Business Ontology (FIBO) provides such a standard in the financial sector [1], it specifies over 70 different types of autonomous agents (entities like persons, organisations, corporations) arranged in a rigorous hierarchy. By defining concepts like \textit{“Legal Entity”} or \textit{“Regulatory Agency”} and their interrelations, FIBO and similar ontologies establish a common vocabulary that banks, insurers, and regulators can share [1,5]. Regulators themselves have advocated adopting common ontological data standards. For example, the Bank of England’s data collection initiative calls for \textit{“defining and adopting common data standards that identify and describe data in a consistent way throughout the financial sector,”} noting that such standards will bring benefits \textit{“well beyond reporting.”} This highlights the industry’s recognition that ontology-backed knowledge graphs can improve interoperability and accuracy in regulatory compliance processes [3,4,6].

        \paragraph{} Knowledge graphs built on ontologies are also used to integrate diverse data sources required for compliance and risk assessment [2,5,6,8]. For instance, \textbf{Know Your Customer} (KYC) and sanctions screening are critical compliance activities in banking [2]. A knowledge graph can link information about people and organisations from internal databases, sanction watch lists, corporate registries, news feeds, and more [5,6]. This creates a unified profile of each entity and exposes connections that might indicate risk. Studies report that knowledge graphs are valuable for such use cases because they provide a \textit{“single access point”} to heterogeneous data and signals [2,6,8]. In a sanctions-checking scenario, a graph-powered system can quickly aggregate data about a client’s beneficial owners, related entities, transaction patterns, and any mentions in adverse media, all in one place. This holistic, connected view enables compliance analysts to detect risks (like a hidden sanctioned party or money laundering network) that would be much harder to spot in isolated data silos [2,6].

        \paragraph{} Ontologies not only unify data but also enable automated reasoning over regulatory knowledge, which is key for scaling compliance checks [3,4,5,6]. A prominent example is \textbf{Automated Compliance Checking} (ACC) in domains like construction and engineering [3]. Here, the goal is to automatically verify whether a project or document complies with a set of regulations. Researchers have shown that encoding regulations in an ontology greatly improves the ability to automate such analysis [3,5]. In one study, regulatory text was parsed into formal semantic triples (subject–predicate–object with conditions) using a domain ontology, yielding a structured representation of each requirement [3]. This ontological model enabled the system to automatically identify duplicate provisions and even detect logical contradictions among requirements. The use of a SKOS-based ontology provided a deeper semantic analysis by capturing hierarchical and associative relationships between terms, going beyond superficial keyword matching [3]. In practical terms, this means the system could recognise when two different wordings impose the same obligation, or when one rule implies the negation of another – tasks that would be extremely difficult with simple text search. By formalising the rules and their terminology, ontologies make it feasible to do consistency checking and highlight compliance issues that human reviewers might overlook, thus mitigating risk [3,4].

        \newpage
        \paragraph{} Combining knowledge graphs with advanced AI further enhances risk and compliance management [4,5,6,11,12]. Large Language Models (LLMs), like GPT-4, are very adept at interpreting and generating text, including legal language, but they can suffer from hallucinations or lack of up-to-date factual knowledge [11,12]. Integrating an LLM with a curated knowledge graph offers the best of both worlds. The knowledge graph provides verified, up-to-date facts and the ontological context, while the LLM provides flexible understanding and explanation in natural language [4,5,11]. For example, an LLM + KG system can answer a query like, \textit{“What new operational risk requirements does regulation X impose on banks?”} by retrieving the exact clauses from the regulatory knowledge graph and then having the LLM present the answer in a concise summary [4,11]. This approach has been shown to improve both accuracy and explainability:\textit{ “Knowledge graphs combined with LLMs address regulatory compliance challenges by automating analysis and providing accurate, explainable guidance that eliminates hallucination risks.”} [4,11] In practice, such a system might be used by a compliance officer or risk manager to quickly get answers with proper citations to the law (since the graph can supply the source). Overall, ontology-backed knowledge graphs, especially when paired with AI language models, are becoming a powerful tool in regulatory compliance – helping institutions stay current with regulations, streamline reporting, identify risks, and ensure that their interpretations of the rules are consistent and well-grounded [4,5,6,11,12].

    \subsection{Digital Twin \& Digital Thread Applications}

        \paragraph{} The \textbf{digital twin} is a virtual mirror of a physical asset or system, maintained in sync with its real-world counterpart’s state. The \textbf{digital thread} refers to the integrated data flow across the life-cycle of a product or system (from design and manufacturing to operation and maintenance). Both concepts rely heavily on integrating and contextualising heterogeneous data – for example, CAD models, simulation results, sensor readings, maintenance logs – into a unified representation. Here, knowledge graphs have emerged as a pivotal enabling technology [13,14,20,21]. They offer a semantic layer to organise and relate all these data elements, essentially threading them together across the product life-cycle. In recent years, knowledge graphs have seen broad application in constructing and managing digital threads [13,14]. By using shared ontologies, a graph can formally represent how a requirement in a design document relates to a component in a CAD model, which in turn relates to a part in an assembly, and further to a sensor measuring that part’s performance in the field [21]. This semantic connectivity is the essence of the digital thread, ensuring that information doesn’t remain in isolated silos but is linked to everything it impacts [20,21].
        
        \newpage
        \paragraph{} A common challenge in implementing digital threads is achieving interoperability among various engineering software and data formats. Ontologies and knowledge graphs address this by serving as an intermediate, standard-based representation [13,14,20]. Researchers have developed methods to convert traditional engineering data schemas into ontologies/graphs for this purpose [13,14]. For example, the \textbf{OntoSTEP} framework takes the STEP CAD data format (an ISO standard for product models) and transforms it into an OWL ontology, thereby creating a knowledge graph of the product’s design. Once in a graph form, it becomes easier to integrate that design data with other sources (such as requirements or manufacturing plans) and to perform queries or inference [21]. Similar efforts have been made for other standards (SysML models, BOMs), all aimed at semantic \textbf{ data integration}. By using ontology-driven graphs, one can ask complex life-cycle questions that span domains – for instance, \textit{“Which design requirements are linked to field issues in the last year?”} – and get an answer because the requirements, design elements, and field reports are all connected in the knowledge graph [13,21]. This approach contrasts with older, siloed IT systems and has been identified as key to a true digital thread [14,20]. In summary, translating heterogeneous data into a common semantic graph allows information to flow with the product across its life-cycle, enabling better change impact analysis, traceability, and systems engineering integration.
        
        \paragraph{} Another major advantage of knowledge graphs in digital twin applications is their ability to handle dynamic, real-time data in a structured way. A physical twin in operation continuously generates new data (sensor measurements, status updates, events), and a digital twin must incorporate these updates. Knowledge graphs support dynamic updates well: new nodes and relationships can be added to represent events or state changes as they occur [14,18,19]. For example, Shen (2025) developed a digital twin of an intelligent structure in which all the knowledge about the asset was managed through an ontology-based graph in neo4j [14]. The graph was not static; as various analysis agents produced results (like a simulation predicting stress in a beam), those outputs were inserted into the knowledge graph with timestamps, effectively logging the twin’s evolving state [18,19]. This timestamped graph serves as a \textbf{living record} of the asset. It enables what is known as a digital thread of the asset’s history – one can trace back through the graph to see what happened at each step in the life-cycle and in what context [14,18]. In general, such dynamic knowledge graphs significantly enhance life-cycle management of assets by preserving the chain of information over time. If an anomaly occurs in the physical twin, engineers can query the graph to see related events or changes that preceded it, all in one interconnected view [14,19]. The ability to fuse temporal events into the knowledge graph and maintain historical context is crucial for diagnostics, maintenance planning, and continuous improvement of the system [25].
        
        \newpage
        \paragraph{} Moreover, knowledge graphs provide a foundation for more intuitive \textbf{human interaction} with digital twin data. Because the graph is inherently semantic, users can query it in high-level terms. There has been exciting progress in combining knowledge graphs with natural language interfaces to query digital twins [12,20,21]. Using an LLM as the front-end and the knowledge graph as the back-end, users can ask questions in plain English and get answers grounded in the twin’s data [24,25]. In the building structure twin mentioned above, the developers implemented a natural language query system where the user’s question is converted (via an LLM-driven tool like LangChain) into a graph query (cypher for neo4j), executed on the knowledge graph, and then the results are narrated by the LLM [20,21]. For example, an engineer could ask, \textit{“Which component showed the highest stress in the last simulation and what were the boundary conditions?”} The system would translate that to a graph query, retrieve the relevant data (e.g. a node representing a component with a \textit{“hasStress”} property and linked to a simulation instance node with certain parameters), and the LLM would generate a succinct answer [12,24]. This retrieval-augmented approach leverages the precise data in the KG and the fluent communication skills of the LLM. Crucially, it ensures the answers are up-to-date and specific to the digital twin, rather than relying on an LLM’s general training data [20,25]. Such integrations hint at the future of digital twin interfaces: interactive, knowledge-driven, and user-friendly, enabling experts to query complex technical data as easily as doing a web search – but with guaranteed relevance and accuracy due to the underlying knowledge graph.
        
    \subsection{System of Systems Applications}

        \paragraph{}
        A \textbf{System of Systems} (SoS) is an integration of independent constituent systems that cooperate for a common goal while retaining their own management and operation. Examples can range from large-scale defence networks, to smart city infrastructures (where transportation, power, water, and communication systems interlink), to multi-enterprise supply chain networks. By their nature, SoS are exceedingly complex: they involve multiple layers of components, various stakeholders, and emergent behaviours that arise from inter-system interactions. Managing an SoS requires understanding not just the parts, but the myriad relationships and dependencies between them. Here, knowledge graphs and related graph technologies have proven invaluable [7,8]. Graph-based data structures and visualisations have emerged as essential tools for handling SoS complexity. They excel at modelling the multi-layered relationships where each component system may impact others in non-obvious ways. For instance, consider a smart city scenario: a major power outage in the electric grid might knock out traffic lights, which in turn affects transportation systems and even emergency services response times. A knowledge graph can explicitly represent these cross-domain links (PowerGrid – feeds → TrafficSystem, TrafficSystem – informs → EmergencyResponse.), allowing analysts to see cascade effects that span different subsystems [13,14]. Graph models are uniquely suited to capture such complex, dynamic relationships, whereas traditional hierarchical or siloed data models might miss them. In essence, the holistic, relationship-centric view provided by a knowledge graph is a natural fit for SoS analysis [7].
        
        \newpage
        \paragraph{}
        Building a knowledge graph for an SoS typically involves unifying data and schemas from very disparate sources – each constituent system might have its own data formats, terminologies, and API's. Ontologies play a key role in harmonising these into a common semantic context [3,5]. By defining overarching concepts and relationships that map to each system’s elements, an ontology-backed knowledge graph can serve as a \textbf{lingua franca} for the SoS. This greatly aids interoperability and data exchange [17]. A recent study by Wu. illustrates this approach: they proposed constructing a \textit{\textbf{“cognitive digital thread”}} for SoS using the Open Services for \textbf{Life-Cycle Collaboration} (OSLC) standard combined with knowledge graph techniques, a vision closely aligned with the digital-thread approaches presented in [14]. The OSLC provides a way to link data across tools (e.g., requirements, architectural models, operational logs), and the knowledge graph provides the semantic integration and reasoning layer on top. The result is that decision-makers can query the integrated knowledge graph to get system-wide insights. For example, one could trace how a requirement in one subsystem (say, a rail network’s safety requirement) propagates constraints or signals into other subsystems (like city emergency planning). This kind of query becomes feasible when all the subsystems’ data are semantically linked. In Wu's case, the approach was shown to support better decision-making and management in a complex SoS by ensuring that relationships and data flows between systems were explicitly captured and could be analysed as a whole. Essentially, the knowledge graph becomes an \textbf{architectural map} of the SoS, one that can be continually updated and queried to understand the system’s behaviour and configuration from a high level down to details [14,17].
        
        \paragraph{}
        Given the scale and complexity of Systems of Systems, there is also interest in leveraging AI to help manage and glean insights from the knowledge graph. One promising direction is using LLMs in conjunction with SoS knowledge graphs to enable more natural queries and automated analyses. An approach known as \textbf{GraphRAG} (graphical Retrieval-Augmented Generation) has been proposed to address SoS complexity, conceptually similar to recent hybrid LLM–KG reasoning models presented in [16,20,23,24,25]. GraphRAG essentially combines large language models with an underlying knowledge graph: the graph stores the trusted knowledge about the SoS (entities, connections, data points), and the LLM provides a flexible interface to that knowledge. In a prototype described by Huang (2025), a complex SoS (Urban Air Mobility network) was managed by constructing a knowledge graph of its components and using the LLM to both interpret user questions and assist in integrating new information into the graph. They found that the knowledge graph could be used to cluster and reduce the apparent complexity (using graph metrics to identify key nodes), and the LLM+KG hybrid outperformed a standard text-only RAG approach in integrating data and answering queries. The implication is that for an SoS, one could ask in plain language, \textit{“What might happen if system A fails while system B is under maintenance?”} and behind the scenes the GraphRAG system would identify relevant links in the knowledge graph (perhaps system A and B share an interface or have a common dependency) and generate a coherent answer grounded in those facts. This kind of AI-driven analysis can augment traditional modelling and simulation by quickly exploring scenarios or pinpointing hidden interdependencies [23,24].
        
        \newpage
        \paragraph{}
        On the human-facing side, graph-based visual analytics are also critical. Visualisation tools help stakeholders navigate the SoS knowledge graph, often through clustering, filtering, or multi-scale views that cut through the complexity [7]. Ultimately, graph visualisation and interaction techniques act as a bridge between the overwhelming complexity of an SoS and human understanding, allowing decision-makers to gain insight and confidence. As one report put it, graph-based visualisation is more than a technical tool — it’s a bridge between overwhelming complexity and human insight, enabling informed decisions even in intricate SoS contexts [8,28]. In summary, knowledge graphs provide the structural backbone to represent Systems of Systems, and when combined with advanced AI for querying and visualisation, they empower us to understand and manage these giant-scale systems in a more transparent and intelligent way [14,23].

    \subsection{Engineering Applications of Knowledge Graphs}
        \paragraph{}
        Engineering practice across domains such as automotive design, aerospace systems, civil infrastructure, industrial manufacturing and energy networks is shaped by an intricate ecosystem of documents, models, standards, simulation outputs and operational observations, and the literature demonstrates that knowledge graphs offer a unifying semantic infrastructure capable of harmonising these sources into a coherent and analytically meaningful representation [13, 14]. By transforming heterogeneous data into structured relational form, knowledge graphs mitigate fragmentation and provide engineers with the ability to trace dependencies between requirements, components, interfaces and historical behaviours with a clarity that traditional document-centric systems rarely support, which is especially valuable in large projects where small oversights can propagate through multiple subsystems before they are detected.

        \paragraph{}
        This integrative capacity is particularly visible in automotive and aerospace industries, where ontologies provide a stable conceptual vocabulary through which requirements, safety constraints, subsystem interactions and regulatory obligations can be expressed consistently across design, verification and compliance teams [13]. Such semantic stability mirrors efforts in the pharmaceutical sector, where regulatory frameworks depend on precise definitions of risk classifications and procedural obligations, and in the construction industry, where building codes must be interpreted uniformly across contractors, engineers and inspection authorities. Across all these contexts the ontology prevents terminological drift and ensures that technical language remains anchored in shared understanding, thereby reducing errors that arise when multiple stakeholders interpret the same documents differently.

        \paragraph{}
        Knowledge graphs further enable the integration of diverse engineering data streams, including simulation results from computational fluid dynamics and finite-element analyses, sensor readings originating from digital twins, lifecycle management records, historical failure logs, maintenance actions and component strategies recorded in enterprise systems [14]. Industries such as aerospace and rail transport rely heavily on this form of structured integration to maintain long-term traceability of component performance, to correlate operational phenomena with design decisions and to align multiple sources of engineering evidence within a single semantic landscape that supports cross-disciplinary understanding.

        \newpage
        \paragraph{}
        Once these rich semantic structures are established, automated reasoning mechanisms become powerful analytical tools, allowing inconsistencies, duplicated constraints, hidden dependencies and structural contradictions to be identified with a degree of precision that is rarely achievable through manual examination alone [14, 21]. In the nuclear energy sector, for example, reasoning over formal safety ontologies can expose latent conflicts between procedural requirements and plant configuration, while in aerospace certification, automated reasoning can highlight missing evidence chains in compliance documentation that might otherwise jeopardise regulatory approval. These analytical capabilities demonstrate that knowledge graphs are not merely passive repositories but active instruments of engineering assurance.

        \paragraph{}
        When knowledge graphs are paired with large language models, the expressive strengths of natural language interpretation combine with the structural reliability of symbolic representation, enabling more articulate synthesis of engineering insights and more grounded explanations of complex design rationales [13, 14, 21]. This hybrid approach is already emerging in semiconductor manufacturing, where LLMs supported by domain graphs summarise failure modes across multiple wafer batches, and in smart-city infrastructure, where graph-grounded LLMs answer queries about interdependent utilities, traffic systems and energy networks while remaining anchored to verifiable domain knowledge. In these contexts, the models offer a balance of fluency and precision that supports not only analysis but also communication across diverse engineering teams.
        
    \subsection{Applied NLP \& LLM Techniques}

        \paragraph{}
        Natural Language Processing and large language model techniques play an essential role in engineering environments where vital domain knowledge resides within standards documents, test reports, interface specifications, contractual requirements and design narratives that are expressed through dense and often specialised textual language [5, 13]. Industries such as pharmaceuticals, aerospace and defence routinely produce thousands of pages of regulatory and procedural documents, and NLP allows these vast corpora to be parsed into entities, relationships and obligations that can be systematically interpreted, integrated and queried within structured knowledge representations.

        \paragraph{}
        Requirement extraction studies show that LLMs are increasingly capable of detecting semantic relationships that are difficult for rule-based systems to capture, such as equivalences hidden within paraphrased text, hierarchical dependencies implicit in multi-clause statements and contextual nuances that reflect domain-specific interpretations of technical language [5, 13]. This capability is crucial in fields like medical device manufacturing, where subtle differences in regulatory wording may correspond to substantial shifts in risk classification, and in aerospace engineering, where functional requirements may be distributed across multiple interdependent documents whose language varies across suppliers, regulators and internal departments. At the same time, scholars warn that without grounding in a domain ontology, LLMs may misinterpret specialised terminology or infer relationships that do not exist, particularly in high-risk domains with evolving regulatory ecosystems [7, 8].

        \newpage
        \paragraph{}
        Through the integration of NLP outputs into knowledge graphs, disparate engineering documents such as interface control documents, validation plans, safety hazard logs, change records and simulation outputs can be unified into a single semantic landscape that reveals interdependencies and historical context across the lifecycle of the system [13, 14]. This consolidation is especially valuable in industries like aerospace maintenance and aviation safety management, where the causes of anomalies often span many layers of documentation and where understanding emerges through the synthesis of information scattered across multiple organisational archives.

        \paragraph{}
        Once textual information is converted into structured knowledge, LLM-supported reasoning can identify contradictions between requirements, reconcile conflicting statements across documents and reveal conceptual links that are not immediately visible through direct human reading [5, 13, 14, 25]. For instance, in pharmaceutical compliance, such systems can highlight discrepancies between manufacturing guidelines and internal operating procedures, while in civil engineering they can detect when building requirements implicitly contradict environmental impact assessments. These reasoning capabilities deepen the model’s contribution by shifting it from a tool for linguistic parsing to an instrument of semantic analysis.

        \paragraph{}
        When tightly integrated with knowledge graphs, LLMs produce explanations that remain anchored in verifiable sources, enabling engineers to query large technical corpora through natural language while maintaining confidence that the system’s responses reflect accurate domain relationships rather than probabilistic speculation [13, 14, 21]. This hybrid configuration is increasingly used in aerospace, smart cities and energy grid management, where domain-grounded LLMs answer complex operational questions with clarity and traceability, thereby strengthening understanding, decision-making and collaboration across interdisciplinary engineering teams.

    \subsection{Applied Reinforcement Learning \& Recommender Systems}
        \paragraph{}
        Reinforcement learning and knowledge-aware recommender systems appear in the engineering literature as adaptive, perceptive companions, particularly suited to domains where decisions unfold not as isolated selections but as sequences interwoven with consequences, constraints and evolving conditions [18, 19, 23]. In additive manufacturing, for example, RL agents observe the delicate interplay between temperature gradients, material deposition patterns and mechanical outcomes, much as a careful observer might study how slight alterations in atmosphere influence the temper of a room. In energy grid management, RL systems follow the shifting ebb and flow of supply, demand and weather patterns, attempting to guide operations toward stability. Meanwhile, in transportation and logistics planning, RL agents respond to congestion, scheduling pressures and resource limitations with strategies that emerge gradually through experience. Recommender systems, when enriched with knowledge graphs, serve a different but complementary purpose, offering suggestions that are rooted in structured domain context, such as identifying which mechanical component is most likely to require inspection or which material configuration best satisfies a set of design constraints [21]. Together, these approaches create a foundation for decision-making in environments that are dynamic, multi-layered and abundant in subtle interdependencies.

        \newpage
        \paragraph{}
        The benefits of these methods are described with considerable admiration in the literature. Reinforcement learning brings an element of discernment that blossoms over time, allowing an agent to refine its behaviour through continual feedback and to discover strategies that may elude explicit programming. In additive manufacturing, such systems can delicately adjust process parameters to improve microstructural quality or reduce defect rates [18, 19]. In railway operations, RL has been used to resolve scheduling conflicts by anticipating how delays propagate across tightly coupled networks. Knowledge-aware recommender systems offer their own charms, for they are not merely pattern matchers but careful scholars of relational structure, capable of drawing upon past configurations, material properties, environmental conditions and design histories to guide engineers toward promising alternatives [21]. Yet these strengths do not arrive without their complications. RL agents can be overly susceptible to noise in their training environment, adopting brittle strategies when feedback is inconsistent or sparse. Recommender systems, if insufficiently grounded in domain knowledge, may prioritise associations that appear statistically appealing while lacking practical or technical relevance. Moreover, both methods must be monitored closely in high-stakes engineering settings where errors carry significant operational, safety or financial consequences.

        \paragraph{}
        Despite these concerns, the literature demonstrates that when reinforcement learning and recommender systems are supported by a carefully constructed knowledge graph, their capacities deepen in a way that feels both elegant and reassuring. In the aerospace sector, hybrid systems can reason about maintenance scheduling by considering not only statistical patterns of wear but also the structural dependencies captured within engineering ontologies. In the energy sector, they can identify how operational adjustments ripple through interconnected subsystems, revealing paths toward optimisation that respect both physical laws and organisational priorities [13, 14]. In manufacturing supply chains they can track how supplier reliability, material behaviour and production timing influence downstream performance and can recommend courses of action that balance efficiency, resilience and cost. By weaving symbolic structure with adaptive learning, these systems produce recommendations that feel thoughtful rather than impulsive, grounded rather than speculative, and articulate rather than opaque.

        \paragraph{}
        Such hybrid approaches enable models to shine a gentle light upon latent dependencies and alternative strategies that would remain concealed in purely numerical analyses. In intelligent transportation planning, they can illuminate how the adjustment of one route reverberates across an entire network of constraints and passenger flows. In predictive maintenance, they can trace patterns of degradation through historical telemetry and environmental exposure, noting which components require attention and which may continue their service unhindered [21, 23, 25]. Their ability to consider multiple steps into the future, evaluating both immediate benefits and long-range implications, gives them a maturity of judgement that classical optimisation methods often lack. Nevertheless, the literature warns that the elegance of these systems must be supported by vigilant oversight. Without careful calibration they may learn to exploit quirks in the data or to amplify subtle biases embedded within historical patterns, particularly in domains where operational decisions have ethical or societal ramifications.

        \newpage
        \paragraph{}
        When these adaptive techniques are paired with knowledge graphs and occasionally complemented by natural language interfaces, they form decision-support systems that are both powerful and pleasingly interpretable. They are increasingly found in smart manufacturing environments, where they coordinate fleets of machines with an appreciation for both efficiency and harmony. They support renewable energy grids, where they balance environmental variability with human demand. They assist aerospace maintenance, offering recommendations that respect the intricate dependencies of aircraft systems. In each of these settings, the fusion of adaptivity, semantic structure and transparent reasoning produces systems that act almost like well-informed advisers, capable of understanding not only what actions are possible but also which are prudent, desirable and aligned with the long-term health of the engineering enterprise [13, 14, 21]. Through this integration, reinforcement learning and recommender systems gain both refinement and restraint, allowing them to operate with a steadiness of judgement befitting the complex, evolving and deeply consequential worlds they inhabit.
        
\section{Integrative Models \& Reference Documents}

    \subsection{Reasoning-Oriented Large Language Models}
        \paragraph{} 
        Reasoning-oriented language models unfold across the literature as a slow crescendo of conceptual refinement, beginning with early knowledge graph research that sought to distil the world into carefully arranged constellations of entities, relations and hierarchical order, thereby turning the fluidity of natural language into a form that could be approached with computational clarity [1]. These foundational ideas established structure not as a mere container of information, but as a cognitive instrument that shapes how both humans and machines perceive meaning, follow dependencies and recognise conceptual boundaries. In the figure that accompanies this discussion, the lowest layer of the diagram represents this structural grounding through its depiction of ontologies, knowledge graphs, regulatory schemas and subsystem models, each portrayed as a steady foundation upon which higher reasoning unfolds. The literature shows how such structure becomes the backbone of more intricate systems that interpret financial and operational risk by linking scattered regulatory obligations into unified semantic landscapes [2], validate compliance by following rule dependencies through logic-based inference [3], or harmonise regulatory contradictions across jurisdictions by making equivalences explicit within a graph [5]. Engineering studies extend this tradition by encoding subsystem obligations, design constraints and functional dependencies in structured forms that prevent reasoning from straying away from domain logic [13]. As these ideas mature, they give rise to neural link-prediction frameworks that infer missing connections across vast relational spaces [16], manufacturing-reasoning systems that read relational patterns from production graphs [18, 19], and hybrid decision-support models that integrate structured representations with data-driven inference for automotive design [21]. Figure 4.1 captures this essential insight visually: the top layers of multi-hop reasoning rest upon the deeper layers of knowledge representation, illustrating the literature’s overarching lesson that meaningful reasoning becomes coherent, deliberate and aligned with domain truth only when it is anchored in explicit, well-crafted knowledge structures [26].

        \begin{figure}[h]
            \centering
            \includegraphics[width=1\linewidth]{Photos/Reasoning LLMs.png}
            \caption{Layered Reasoning Architecture Over A Knowledge Graph}
            \label{fig:reason}
        \end{figure}

        \newpage
        \paragraph{} 
        From this body of work we learn that the presence of structure grants reasoning-oriented models a stability and conceptual discipline that purely neural systems often struggle to maintain, particularly when tasks require the delicate choreography of multi-step inference. In the middle layer of the diagram, we illustrate this interplay through modules responsible for semantic parsing, dependency recognition, graph-grounded inference and cross-document mapping, each of which operates by returning continuously to the structured foundation beneath. The literature shows that when models are guided by ontologies or graph schemas they move with a clarity that emerges from knowing which concepts exist, how they are related and what constraints govern their use, which is especially important in regulatory analysis [3, 4], requirement harmonisation [5] and cross-domain integration [7, 8] where each inferential step must be interpretable and traceable. Studies demonstrate that knowledge graphs enable a model to reveal its reasoning in identifiable steps rather than in a haze of opaque neural intuition, which gives auditors and engineers the ability to understand how conclusions were reached [13, 18]. Yet the literature equally reminds us that neural reasoning retains vulnerabilities that grow sharper as inference chains lengthen. In additive manufacturing, models may appear coherent while producing predictions that collapse under scrutiny when underlying data are incomplete or ambiguous [18, 19]. In policing analytics, gaps in the graph can lead to interpretations with dangerous real-world consequences [28]. In engineering design evaluation, models may cling to patterns that appear meaningful but fail under rigorous examination [21]. These cautions reinforce what the figure attempts to visually communicate: reasoning requires not only a structured foundation, but also continuous validation mechanisms, checks against hallucination and the capacity to detect when a reasoning chain has drifted into conceptual implausibility, particularly in safety-critical contexts [26].

        \newpage
        \paragraph{}
        As the field evolves, reasoning-oriented models begin to take on an increasingly agentic and exploratory character, behaving less like static predictors and more like travellers moving through carefully built landscapes of structured meaning. The topmost layer of the figure illustrates this through the depiction of a flowing reasoning path that begins with an extracted ADR clause and moves gracefully through stages such as regulatory interpretation, subsystem obligation identification, design-dependency mapping, cross-rule integration, risk or compliance checks and finally the emergence of a reasoned output. The literature shows how models equipped with relational embeddings and graph-informed representations move through rulebooks with an awareness of dependencies and exceptions that previously required manual interpretation [5], how they map multi-modal engineering ecosystems by recognising the interplay of mechanical, electrical and software constraints [13], how they follow temporal sequences within digital twins while maintaining coherence between past and present [14], and how they infer manufacturing behaviour from relational signatures woven into process graphs [18, 19]. Studies reveal that such models achieve improved accuracy, consistency and interpretability in multi-hop reasoning tasks when they are grounded in explicit structure [16, 25]. The broader lesson, both from the literature and from the figure itself, is that reasoning becomes most powerful when the model inhabits the structure it is given, respecting its boundaries while drawing upon neural generalisation to fill in what is implicit. These models move not as improvisers in an empty space, but as attentive participants within a carefully shaped conceptual world, blending symbolic clarity with neural expressiveness to produce reasoning that is ultimately more deliberate, more transparent and more trustworthy for domain-specific intelligence.

    \newpage
    \subsection{Structure \& Planning Oriented Large Language Models}

        \paragraph{} 
        Structure and planning oriented language models emerge in the literature as deliberate responses to the layered intricacy of engineering and organisational environments, where tasks unfold as interconnected sequences of decisions rather than isolated prompts. These works show that structured compliance workflows are strengthened when obligations, constraints and allowable conditions are formalised in shared ontologies that allow models to follow logical instruction chains [4]. Regulatory harmonisation studies use structured vocabularies to bridge discrepancies between rulesets, enabling models to follow equivalence mappings with precision [5]. In automotive engineering, planning becomes possible only when subsystem responsibilities, design constraints and requirement hierarchies are encoded into a consistent semantic structure that the model can reference as it navigates decisions [13]. Digital thread literature demonstrates that lifecycle planning relies on semantically integrated data, allowing the model to see how choices made early in design propagate into manufacturing, testing and operation [14]. Meanwhile, KG supported design recommendation frameworks show that structured pathways allow the model to sequence suggestions based on technical feasibility, historical patterns and cross component dependencies [21]. Figure 4.2 visually echoes this body of work by illustrating how a planning query enters a staged architecture, travelling first through structured substrates such as regulatory ontologies, knowledge graphs, subsystem models and digital thread structures, before being processed by directed reasoning modules that perform constraint retrieval, hierarchy navigation, dependency mapping, temporal reasoning and cross document alignment. The diagram embodies the very principles identified in the literature by showing that planning intelligence arises from moving systematically across layers of structured knowledge rather than improvising in isolation. Together these papers teach us that planning capacity emerges naturally when a model is not left to improvise freely but is supported by domain structure that shapes the flow of its reasoning.

        \begin{figure}[h]
            \centering
            \includegraphics[width=0.95\linewidth]{Photos/Structured Planning Pipeline.png}
            \caption{Structured Planning Pipeline For Planning-Oriented LLMs}
            \label{fig:structure}
        \end{figure}

        \paragraph{}
        The literature highlights that structured alignment yields significant advantages, especially in high stakes sectors where planning errors can lead to financial loss, safety incidents or systemic failures. Automated compliance interpretation benefits from structured legal hierarchies that prevent the model from misreading conditional statements [4]. Requirement mapping becomes clearer and more scalable when guided by ontologies that encode domain relationships [5]. Engineering coordination becomes more accurate when subsystems are linked through explicit graph pathways rather than informal documentation [13]. However, the papers also illuminate the limitations of overly rigid structural frameworks. Regulatory studies show that schemas can quickly become outdated as rules evolve, making it difficult for models to adapt unless the structure itself is dynamically updated [7, 8]. Interactive planning systems demonstrate that models can struggle when user intent conflicts with the strict logic encoded in the structure [21]. Manufacturing research highlights the difficulty of capturing real operational variability within static graphs or ontologies [18, 19]. Figure 4.2 further reinforces these insights by depicting how the transition from structured substrates to neural planning modules requires a subtle balance between rigidity and adaptability. The arrows that bridge the layers emphasise that structure alone cannot sustain planning performance unless the system retains the flexibility to reinterpret, refine and negotiate those structures as conditions shift, thereby echoing the literature’s call for frameworks that combine stability with creative adaptability. From these insights the literature teaches us that structure must remain flexible and adaptive, allowing neural reasoning to fill gaps and adjust to unseen scenarios while still maintaining the interpretability and traceability that structured frameworks provide.
        
        \paragraph{}
        Across domains structure oriented language models are presented as quiet orchestrators of multi stage reasoning, weaving together knowledge from diverse parts of a system and sequencing it into coherent, well grounded plans. Studies show how models map design dependencies across entire engineering lifecycles by following graph based structures that encode component relationships [13]. Regulatory mapping frameworks demonstrate how models can follow chains of obligations linking multiple documents and rule sources [5]. Digital thread systems highlight how structured representations allow models to plan across temporal stages of a product, understanding how design decisions influence later operational behaviour [14]. Other works show how structural reasoning supports simulation based analysis [22], hybrid symbolic neural methods that integrate domain rules with neural inference [20], and reinforcement learning systems that rely on structured states to guide efficient exploration [23]. Figure 4.2 crystallises these ideas by presenting planning not as a single operation but as a sequenced progression, where each reasoning step draws nourishment from the structured knowledge below it and contributes to a refined, defensible plan at the output. The figure therefore stands as a conceptual mirror to the literature, illustrating that planning oriented LLMs achieve their elegance and reliability precisely by traversing the carefully designed knowledge landscape that underlies each decision. The broader lesson from these studies is that planning becomes most elegant and reliable when the model is not navigating an empty landscape but instead traversing a carefully designed knowledge structure that provides orientation, stability and conceptual depth.

    \subsection{Memory \& Agent-Based Large Language Models}

        \paragraph{} 
        Memory-augmented and agent-based language models appear in the literature as attempts to inject continuity, history and long-term intention into the reasoning process, acknowledging that genuine intelligence cannot emerge if the model forgets the world at the end of every prompt. Research on autonomous memory retrieval highlights how models benefit from mechanisms that allow them to re-access past states and reasoning traces, thereby reducing conceptual fragmentation and improving consistency [11]. Studies on conversational persistence illustrate how continuity across interactions leads to more stable, context-aware interpretation of user needs [24]. Agentic frameworks demonstrate that when models are given autonomy to refine their own strategies, they develop more coherent reasoning patterns that evolve over time [12]. KG-anchored memory architectures show that storing relational and temporal context helps models reason about not only static facts but also how those facts relate to unfolding events [25]. These insights are echoed across regulatory assessment, where memory allows the model to track cross-document obligations [4]; engineering requirement management, where continuity helps maintain coherence across design iterations [13]; digital-twin modelling, where temporal context is inseparable from system interpretation [14]; and manufacturing-process analysis, where understanding requires remembering multi-stage dependencies [18, 19]. 

        \paragraph{} 
        A particularly illustrative representation of these ideas is captured in Figure 4.3, which gathers the themes of the literature into a single architectural landscape and renders visible the way memory and agency intertwine to support coherent long-horizon reasoning. The figure presents a concentric structure beginning with a central memory core that holds entity histories, temporal records, validated facts, prior reasoning chains and contradiction flags, and this nucleus is encircled by modular agents such as retrieval, reasoning, planning, consistency checking and domain validation, each circulating information in a continuous rhythm of querying, updating and aligning. Surrounding these agents sits an outer band of real-world contexts in which memory is indispensable, including regulatory assessment [4], engineering requirement management [13], digital-twin evolution [14], multi-stage manufacturing processes [18, 19], conversational persistence [24] and agentic behaviour refinement [12], reminding us that the value of internal continuity becomes fully apparent only when the model must navigate environments where past states govern present interpretation. This visual architecture mirrors what the literature repeatedly reveals, which is that these systems behave most intelligently when memory is not treated as a transient buffer but as a living substrate that shapes how agents perceive, reason and act across time, allowing them to weave together temporal cues, relational signals and previously validated decisions into a more stable and reflective form of computational thought.
        
        \begin{figure}[h]
            \centering
            \includegraphics[width=1\linewidth]{Photos/Memory & Agent Based LLMs.png}
            \caption{Multi-Agent Memory \& Temporal Reasoning Architecture}
            \label{fig:memory}
        \end{figure}

        \newpage
        \paragraph{} 
        The literature highlights that persistent memory offers substantial advantages, improving multi-hop reasoning accuracy by enabling the model to reference earlier steps rather than reconstruct context from scratch [25]. Structured memory retrieval mitigates hallucination by grounding new claims in previously verified information [11]. Agent-based planning becomes more stable when the model maintains an internal state that shapes its evolving behaviour [12, 24]. Yet the papers also warn that memory introduces new vulnerabilities. Incorrect or outdated information can accumulate silently if memory is not regularly validated, creating distortions that weaken later reasoning. This is particularly dangerous in policing analytics, where incorrect inferences can lead to harmful real-world actions [28], in manufacturing reasoning where a single incorrect dependency can propagate through entire workflows [18, 19], and in regulatory interpretation where accuracy is non-negotiable [4]. Therefore, the literature teaches us that memory must be continuously curated, updated, checked for contradictions and aligned with domain truth to prevent error accumulation and conceptual drift.

        \paragraph{} 
        Across all domains memory-augmented and agent-based systems represent not merely enhancements but fundamental shifts in the way language models reason, enabling them to sustain coherent thought across time, integrate new insights with past learning and respond to complex tasks with accumulated wisdom. Their influence can be seen in ontology-driven construction workflows that rely on long-term context preservation to ensure consistency across design stages [17], in digital-twin architectures that retain event histories and use them to interpret current system behaviour [14], in engineering requirement systems that depend on stable anchors to maintain coherence across subsystem interactions [13], in manufacturing analysis models that adjust to variable conditions through learned state information [18, 19], and in design-recommendation approaches where memory enables richer, context-dependent suggestions [21]. Collectively these papers suggest that memory and agency are guiding language models toward a form of sustained computational intelligence that moves fluidly across time, participates in iterative reasoning cycles and inhabits its knowledge environment with increasing nuance, adaptability and conceptual depth.
        
    \subsection{Hy-Brid KG–Neural Network Architectures}

        \paragraph{} 
        Hybrid knowledge graph and neural network architectures emerge throughout the literature as an attempt to braid together two fundamentally different ways of representing knowledge, combining the clarity and explicitness of symbolic structures with the fluid, pattern-seeking capabilities of deep learning. The papers that examine KG embeddings demonstrate how neural encoders can learn dense vector representations of entities and relations that preserve graph structure while enabling efficient computation over large and complex domains [16]. Studies on predictive manufacturing and engineering analytics show how these embeddings allow neural models to capture subtle relational patterns, such as latent dependencies in additive manufacturing processes or correlations between geometric features and defect likelihoods [18, 19]. Other works highlight the power of relational graph convolutional networks in modelling cross-domain interactions within engineering systems, illustrating that when the graph’s relational form is given a neural expressive layer, the resulting architecture can reason over both explicit structure and implicit statistical signatures with improved accuracy [25]. Across these studies a central insight emerges: hybrid architectures are not meant to replace symbolic or neural systems, but to elevate both by giving symbolic knowledge the flexibility of learning and giving neural learning the discipline of structure.

        \newpage
        \begin{figure}[h]
            \centering
            \includegraphics[width=0.75\linewidth]{Photos/Hy-Brid KG-N.png}
            \caption{Hy-Brid Neuro-Symbolic Reasoning Pipeline}
            \label{fig:Hy-Brid Neuro-Symbolic Reasoning Pipeline}
        \end{figure}

        \paragraph{} 
        From these papers we learned that hybrid approaches offer substantial benefits, particularly in domains where explicit structure exists but raw data also contains hidden patterns that humans cannot easily formalise. For compliance, hybrid models can use the symbolic layer to ensure factual accuracy while using the neural layer to detect contextual signals, ambiguity, or subtle linguistic cues in regulatory documents that might indicate exceptions or edge cases [3, 4]. In engineering, hybrid reasoning captures both the deterministic dependencies encoded in design ontologies and the non-linear behaviours observed in real-world sensor data or manufacturing logs [13, 18, 19]. This dual-layered reasoning improves interpretability, stability and predictive power. However, the literature also exposes notable risks. Neural embeddings can distort or oversimplify graph semantics if trained without constraint, creating representations that capture statistical regularities yet lose conceptual meaning [16]. In dynamic environments, hybrid models may inherit fragility from both sides: symbolic structures may become outdated, while neural components may overfit to transient patterns. Studies in regulatory and policing contexts warn that hybrid systems can amplify biases if the graph encodes incomplete or skewed knowledge or if neural layers generalise from biased historical data [28]. The overarching lesson is that hybrid KG–neural architectures must be carefully balanced, continuously updated and monitored through robust alignment mechanisms to prevent conceptual drift, representational collapse or unintended amplification of errors.

        \newpage
        \paragraph{} 
        Across the literature hybrid architectures are presented as precursors to a more tightly integrated neuro-symbolic paradigm, where symbolic scaffolding and neural adaptability work together not as separate modules but as co-evolving components of a unified reasoning system. Engineering-focused studies envision hybrid models capable of simulating entire system behaviours by combining ontology-grounded constraints with deep neural predictions, enabling models to understand not only the structure of a subsystem but also how its behaviour changes under different conditions [13, 14]. Research in design recommendation frameworks shows that hybrid systems can generate contextually meaningful suggestions by following graph-encoded best practices while using neural priors to identify promising alternatives [21]. Digital-twin literature suggests that hybrid architectures can synthesise symbolic state representations with time-sensitive neural features to model evolving system dynamics [14]. Meanwhile, work on advanced link-prediction and multi-hop reasoning demonstrates that hybrid systems improve inference accuracy when navigating incomplete or evolving knowledge spaces [16, 25]. Taken together, these papers reveal a trajectory in which hybrid architectures become engines of adaptive reasoning, capable of integrating explicit knowledge, learned behaviour and temporal evolution into a single coherent system. They hint at an emerging vision where intelligent systems no longer oscillate between symbolic rigidity and neural intuition, but instead learn to inhabit the shared space between them with increasing sophistication, nuance and conceptual depth.

        